\subsection*{導入：モデルと方法が重要な理由}
\addcontentsline{toc}{subsection}{導入：モデルと方法が重要な理由}
ソフトウェア開発では、多くの人が同時に異なる観点で考える。利用者は業務の流れを考え、設計者は構造を考え、開発者はコードの形を考え、テスト担当者は確認方法を考える。これらが頭の中だけにあると、誤解や抜け漏れが起きやすい。

モデルは、この問題に対する強力な道具である。モデルは対象の全体を完全に写すのではなく、目的に合わせて重要な側面を選び取る。たとえば、注文管理システムを考えるとき、データ構造を考えるなら「注文」「商品」「顧客」の関係を図にする。処理の流れを考えるなら「注文受付」「在庫確認」「支払い」「出荷」の順序を図にする。どちらも同じシステムのモデルだが、見ている側面が違う。

方法は、モデルや成果物をどの順序で作り、どのように確認し、どう次の作業へ渡すかを決める枠組みである。経験則に基づく方法、形式手法、プロトタイピング、アジャイル方法は、いずれも「よりよく作る」ための方法だが、適している状況が異なる。

\begin{pointbox}{重要}
この章でいうモデルは、きれいな図を描くこと自体が目的ではない。モデルの目的は、考える、説明する、確認する、合意する、変更の影響を追うことである。使われないモデルは価値が低い。
\end{pointbox}
